\documentclass[11pt]{article}
\usepackage[margin=1in]{geometry}
\usepackage[T1]{fontenc}
\usepackage[utf8]{inputenc}
\usepackage{lmodern}
\usepackage{microtype}
\usepackage{amsmath,amssymb,amsthm,mathtools}
\usepackage{enumitem}
\usepackage[hidelinks]{hyperref}
\usepackage{xcolor}
\usepackage{listings}

\lstdefinestyle{lean}{
  basicstyle=\ttfamily\small,
  breaklines=true,
  columns=fullflexible,
  keepspaces=true,
  showstringspaces=false,
  frame=single
}

\theoremstyle{plain}
\newtheorem{theorem}{Theorem}
\newtheorem{remark}{Remark}

\newcommand{\RR}{\mathbb{R}}
\newcommand{\Rp}{\RR_{>0}}
\newcommand{\Jcost}{J}

\title{Lean Support for the Full Unconditional Closure Theorem}
\author{}
\date{\today}

\begin{document}
\maketitle
\emergencystretch=3em

\section*{Purpose}
This note records the Lean framework in the repository that supports the
``full unconditional closure'' result stated in
\path{RS_CostAlgebra_Recognition_Category.tex}, and it gives a complete
written proof in ordinary mathematical language.

\section*{Supporting Lean Modules}
The relevant public modules are:
\begin{itemize}[nosep]
\item \path{RecognitionScience/Cost/FunctionalEquation.lean}
\item \path{RecognitionScience/Foundation/DAlembert/Unconditional.lean}
\item \path{RecognitionScience/Foundation/DAlembert/FullUnconditional.lean}
\item \path{RecognitionScience/Cost/AczelTheorem.lean}
\end{itemize}

The key theorem names are:
\begin{itemize}[nosep]
\item \path{dAlembert_cosh_solution_aczel}
\item \path{washburn_zlatanovic_uniqueness_aczel}
\item \path{J_computes_P}
\item \path{J_surjective_nonneg}
\item \path{P_determined_nonneg}
\item \path{consistency_forces_RCL_form_is_theorem}
\item \path{full_unconditional_inevitability}
\item \path{washburn_zlatanovic_full_unconditional}
\end{itemize}

\section*{Lean Imports}
\begin{lstlisting}[style=lean]
import RecognitionScience.Cost.FunctionalEquation
import RecognitionScience.Foundation.DAlembert.Unconditional
import RecognitionScience.Foundation.DAlembert.FullUnconditional
\end{lstlisting}

\section*{Main Lean Theorem}
The public theorem in the repository closest to the paper statement is
\path{washburn_zlatanovic_full_unconditional}. For portability in this
document, the statement below is written in ASCII Lean style while preserving
the theorem name and hypotheses.

\begin{lstlisting}[style=lean]
theorem washburn_zlatanovic_full_unconditional
    (F : Real -> Real)
    (P : Real -> Real -> Real)
    (hSymm : forall x : Real, 0 < x -> F x = F (1 / x))
    (hUnit : F 1 = 0)
    (hSmooth : ContDiff Real 2 F)
    (hCalib : deriv (deriv (G F)) 0 = 1)
    (hCons : forall x y : Real, 0 < x -> 0 < y ->
      F (x * y) + F (x / y) = P (F x) (F y))
    (hPoly : exists a b c d e f : Real,
      forall u v, P u v = a + b*u + c*v + d*u*v + e*u^2 + f*v^2)
    (hSymP : forall u v, P u v = P v u)
    (hNonTriv : exists x : Real, 0 < x /\ F x /= 0)
    (hCont : ContinuousOn F (Set.Ioi 0))
    (hP11 : P 1 1 = 6) :
    (forall x : Real, 0 < x -> F x = Cost.Jcost x) /\
    (forall u v : Real, 0 <= u -> 0 <= v ->
      P u v = 2 * u * v + 2 * u + 2 * v)
\end{lstlisting}

Its body is the short wrapper:

\begin{lstlisting}[style=lean]
refine full_inevitability_explicit F P hSymm hUnit hSmooth hCalib hCons ?_ ?_
  exact consistency_forces_RCL_form_is_theorem F P hSymm hUnit hSmooth
    hCons hPoly hSymP hNonTriv hCont hP11
  exact dAlembert_forces_cosh_is_theorem
\end{lstlisting}

\section*{Auxiliary Lean Results}
The two structural ingredients used by the wrapper above are:

\begin{lstlisting}[style=lean]
theorem consistency_forces_RCL_form_is_theorem
    (F : Real -> Real) (P : Real -> Real -> Real)
    (_hSymm : forall x : Real, 0 < x -> F x = F (1 / x))
    (hUnit : F 1 = 0)
    (_hSmooth : ContDiff Real 2 F)
    (hCons : forall x y : Real, 0 < x -> 0 < y ->
      F (x * y) + F (x / y) = P (F x) (F y))
    (hPoly : exists a b c d e f : Real,
      forall u v, P u v = a + b*u + c*v + d*u*v + e*u^2 + f*v^2)
    (hSymP : forall u v, P u v = P v u)
    (hNonTriv : exists x : Real, 0 < x /\ F x /= 0)
    (hCont : ContinuousOn F (Set.Ioi 0))
    (hP11 : P 1 1 = 6) :
    forall x y : Real, 0 < x -> 0 < y ->
      P (F x) (F y) = 2 * F x * F y + 2 * F x + 2 * F y
\end{lstlisting}

\begin{lstlisting}[style=lean]
theorem dAlembert_cosh_solution_aczel
    (H : Real -> Real)
    (h_one : H 0 = 1)
    (h_cont : Continuous H)
    (h_dAlembert : forall t u, H (t + u) + H (t - u) = 2 * H t * H u)
    (h_d2_zero : deriv (deriv H) 0 = 1) :
    forall t, H t = Real.cosh t
\end{lstlisting}

The unconditional determination of the combiner once the canonical cost is known
is handled in:

\begin{lstlisting}[style=lean]
theorem J_computes_P :
    forall x y : Real, 0 < x -> 0 < y ->
      Cost.Jcost (x * y) + Cost.Jcost (x / y) =
      2 * Cost.Jcost x * Cost.Jcost y +
      2 * Cost.Jcost x + 2 * Cost.Jcost y

theorem J_surjective_nonneg :
    forall v : Real, 0 <= v ->
      exists x : Real, 0 < x /\ Cost.Jcost x = v

theorem P_determined_nonneg (P : Real -> Real -> Real)
    (hCons : forall x y : Real, 0 < x -> 0 < y ->
      Cost.Jcost (x * y) + Cost.Jcost (x / y) =
      P (Cost.Jcost x) (Cost.Jcost y)) :
    forall u v : Real, 0 <= u -> 0 <= v ->
      P u v = 2*u*v + 2*u + 2*v
\end{lstlisting}

\section*{Mathematical Statement}
\begin{theorem}[Full unconditional closure]
Let $F:\Rp \to \RR$ be $C^2$, and let $P:\RR^2 \to \RR$ be symmetric.
Assume:
\begin{enumerate}[label=\textup{(\roman*)},nosep]
\item $F(1)=0$,
\item $F(x)=F(x^{-1})$ for all $x>0$,
\item $\left.\dfrac{d^2}{dt^2}\right|_{t=0} F(e^t)=1$,
\item $F(xy)+F(x/y)=P(F(x),F(y))$ for all $x,y>0$,
\item $F$ is nontrivial: $F(x_0)\neq 0$ for some $x_0>0$,
\item $P(1,1)=6$.
\end{enumerate}
If, in addition, $P$ lies in the quadratic symmetric ansatz used by the public
Lean theorem,
\[
P(u,v)=a+b\,u+c\,v+d\,uv+e\,u^2+f\,v^2,
\]
then
\[
F(x)=\frac12(x+x^{-1})-1 \qquad (x>0),
\]
and
\[
P(u,v)=2uv+2u+2v \qquad (u,v\ge 0).
\]
\end{theorem}

\begin{remark}
The public Lean development visible in the repository proves the theorem above.
The d'Alembert-to-cosh step is packaged through
\path{aczel_representation_axiom} in
\path{RecognitionScience/Cost/AczelTheorem.lean}, and the combiner-forcing
step in \path{consistency_forces_RCL_form_is_theorem} uses an internal
bilinear-forcing result referenced in
\path{RecognitionScience/Foundation/DAlembert/FullUnconditional.lean}.
\end{remark}

\section*{Full Written Proof}
\begin{proof}
Define
\[
G(t):=F(e^t) \qquad (t\in\RR).
\]
Then $G$ is $C^2$, because $F$ is $C^2$ on $(0,\infty)$ and $t\mapsto e^t$ is
smooth. The normalization gives
\[
G(0)=F(1)=0.
\]
The reciprocal symmetry of $F$ implies that $G$ is even:
\[
G(-t)=F(e^{-t})=F((e^t)^{-1})=F(e^t)=G(t).
\]
Hence $G'(0)=0$, and the calibration hypothesis becomes
\[
G''(0)=1.
\]

Because $G$ is $C^2$ and satisfies
\[
G(0)=0,\qquad G'(0)=0,\qquad G''(0)=1,
\]
its second-order Taylor expansion at $0$ is
\[
G(t)=\frac{t^2}{2}+o(t^2)\qquad (t\to 0).
\]
Therefore $G(t)>0$ for all sufficiently small nonzero $t$. Since $G$ is
continuous and $G(0)=0$, it follows that the range of $G$ contains a nontrivial
interval $[0,\varepsilon]$. Equivalently, the range of $F$ on $(0,\infty)$
contains a nontrivial interval $[0,\varepsilon]$.

Now use the quadratic symmetric ansatz from the Lean theorem. Because $P$ is
symmetric, it can be written in the form
\[
P(u,v)=a+b(u+v)+cuv+d(u^2+v^2)
\]
for suitable real constants $a,b,c,d$.

Set $y=1$ in the consistency law. Since $F(1)=0$, we obtain
\[
F(x)+F(x)=P(F(x),0),
\]
that is,
\[
2F(x)=P(F(x),0)\qquad (x>0).
\]
Because the range of $F$ contains an interval $[0,\varepsilon]$, this implies
\[
P(u,0)=2u\qquad \text{for all }u\in[0,\varepsilon].
\]
But
\[
P(u,0)=a+bu+du^2.
\]
Hence the polynomial identity
\[
a+bu+du^2=2u
\]
holds on an interval. Therefore the two polynomials are identical, so
\[
a=0,\qquad b=2,\qquad d=0.
\]
Thus $P$ simplifies to
\[
P(u,v)=2u+2v+cuv.
\]

Now evaluate at $(u,v)=(1,1)$. The normalization $P(1,1)=6$ gives
\[
6=P(1,1)=2+2+c,
\]
so $c=2$. Consequently,
\[
P(u,v)=2uv+2u+2v.
\]
Substituting this into the consistency law yields
\[
F(xy)+F(x/y)=2F(x)F(y)+2F(x)+2F(y)\qquad (x,y>0),
\]
which is exactly the Recognition Composition Law.

At this point the standard uniqueness theorem applies. The Lean theorem
\path{washburn_zlatanovic_uniqueness_aczel} in
\path{RecognitionScience/Cost/FunctionalEquation.lean} proves that any
continuous reciprocal cost satisfying the Recognition Composition Law,
normalization, and calibration must equal the canonical cost
\[
\Jcost(x)=\frac12(x+x^{-1})-1.
\]
Since $F$ is $C^2$, it is continuous, so we conclude
\[
F(x)=\frac12(x+x^{-1})-1 \qquad (x>0).
\]

It remains to show that the same formula for $P$ holds on the entire first
quadrant $[0,\infty)^2$. Fix $u,v\ge 0$, and define
\[
x_u:=(u+1)+\sqrt{(u+1)^2-1},
\qquad
y_v:=(v+1)+\sqrt{(v+1)^2-1}.
\]
These numbers are strictly positive. A direct calculation shows that
\[
\Jcost(x_u)=u,\qquad \Jcost(y_v)=v,
\]
because solving
\[
\frac12(x+x^{-1})-1=u
\]
for $x>0$ gives exactly $x=x_u$, and similarly for $v$.

Therefore
\[
P(u,v)=P(\Jcost(x_u),\Jcost(y_v)).
\]
Using the consistency law with $F=\Jcost$, we obtain
\[
P(\Jcost(x_u),\Jcost(y_v))
=\Jcost(x_u y_v)+\Jcost(x_u/y_v).
\]
Now apply the Recognition Composition Law for $\Jcost$:
\[
\Jcost(x_u y_v)+\Jcost(x_u/y_v)
=2\Jcost(x_u)\Jcost(y_v)+2\Jcost(x_u)+2\Jcost(y_v).
\]
Since $\Jcost(x_u)=u$ and $\Jcost(y_v)=v$, this becomes
\[
P(u,v)=2uv+2u+2v.
\]
Thus both the scalar law and the combiner are uniquely determined.
\end{proof}

\section*{Summary}
The supporting Lean chain is:
\begin{enumerate}[nosep]
\item the combiner on the range of $F$ is forced to the RCL form by
\path{consistency_forces_RCL_form_is_theorem},
\item the scalar law is forced to the canonical cost by
\path{dAlembert_cosh_solution_aczel} and
\path{washburn_zlatanovic_uniqueness_aczel},
\item the combiner is then forced on all of $[0,\infty)^2$ by
\path{J_surjective_nonneg} and \path{P_determined_nonneg}.
\end{enumerate}

\end{document}
